% !TeX root = ../tesis.tex

\section{Conclusiones}
\label{sec:conclusions-and-future-work-conclusiones}

La conclusión principal de este trabajo es que sí es posible materializar un \zkBridge sobre \Cardano respetando los límites de unidades de ejecución impuestos por el protocolo. En particular, las mediciones del Capítulo~\ref{chap:resultados} muestran que, mediante la partición de la verificación \STM en dos fases de la Subsección~\ref{subsec:mithril-a-cardano-particion-de-la-verificacion-en-dos-transacciones}, el uso de \ReferenceScripts y la separación entre capa de actualización y capa de aplicación de la Subsección~\ref{subsec:design-decisions-capas-actualizacion-y-aplicacion}, el flujo final puede mantenerse dentro del presupuesto \Onchain disponible sin abandonar el modelo \eUTxO.

\subsection{Evaluación del modelo de seguridad y alcance del prototipo}
\label{subsec:conclusions-and-future-work-modelo-de-seguridad-y-alcance-del-prototipo}

El modelo de amenaza, sus supuestos y las propiedades buscadas se fijaron antes de presentar la arquitectura, en la Sección~\ref{sec:design-decisions-modelo-de-seguridad}. Bajo ese modelo, la implementación aporta evidencia concreta para las cinco propiedades allí enumeradas. La autenticidad del estado de certificación se obtiene mediante una \CertificateChain anclada en un \GenesisCertificate y continuada por \StakeDistributionUTxO sucesivos. El binding entre verificación y consumo se materializa mediante \ProofReceipt producidos por la verificación \STM en dos fases: \PhaseOneSetupTx fija el estado intermedio del \transcript, y \PhaseTwoVerifyTx sólo produce el receipt si el \ReducedRedeemer ligado a ese estado completa satisfactoriamente el chequeo final. El \texttt{statement\_hash}, el dominio y la política del receipt conectan ese resultado con el mensaje del certificado que consumirá la transacción posterior.

Para el camino $\LockTx \rightarrow \MintTx$, cuyo ensamblado completo se presenta en la Subsección~\ref{subsec:final-bridge-flow-el-flujo-principal-del-zkbridge}, el argumento de pertenencia exige que el \TxId reclamado aparezca en un \TransactionSnapshot autenticado por \Mithril. La interpretación de ese identificador como una operación concreta de \textit{locking} se completa mediante la interfaz de los validadores: el reclamo de aplicación vincula el identificador, el valor y el destinatario que la formulación del prototipo puede reconstruir. Por último, \MintTx consume el \TxsUpdaterUTxO vigente y verifica la transición del \SMT de ausencia a presencia del mismo \TxId. Como el minting y esa actualización ocurren dentro de una única transacción \eUTxO, la continuidad del estado y la protección contra repeticiones se aceptan o rechazan conjuntamente.

Estas garantías son condicionales a los supuestos de la Tabla~\ref{tab:bridge-security-assumptions}: autenticidad de \GenesisVerificationKey, seguridad de \Mithril bajo su modelo de stake, solidez de los argumentos \HaloTwo/\Groth, resistencia de las funciones de hash, codificaciones canónicas y corrección de circuitos, claves y validadores desplegados. No se supone que el operador, el \textit{relayer}, el \textit{prover} o el \MithrilAggregator sean honestos para interpretar los datos que presentan; los scripts deben rechazar evidencia inválida. Sí se necesita que exista infraestructura disponible para producir argumentos y enviar transacciones. Por ello, el prototipo separa seguridad de \textit{liveness}: un actor \Offchain puede detener o censurar el progreso, pero su indisponibilidad no constituye por sí sola una forma de producir un minting aceptado sin evidencia.

El alcance del prototipo, por lo tanto, debe leerse como una prueba de viabilidad para una instancia de \Cardano, no como un bridge productivo y heterogéneo listo para cualquier par de cadenas. La implementación trabaja con una formulación isomorfa de las transacciones relevantes: las cadenas $\ChainA$ y $\ChainB$ pueden modelarse con estructuras suficientemente compatibles para reconstruir y validar localmente el evento de origen. En un escenario plenamente heterogéneo, esa capa debería reemplazarse por circuitos y codificaciones específicos de la cadena origen, junto con reglas explícitas para normalizar direcciones, valores, fees y formatos de transacción.

En síntesis, el prototipo demuestra, bajo los supuestos declarados y para el flujo unidireccional implementado, que es posible componer \Onchain certificados \Mithril, recibos de verificación y argumentos sucintos dentro de los límites reales de \Cardano. Las extensiones necesarias para convertir esta prueba de concepto en un sistema operativo sobre una red pública se desarrollan en la Sección~\ref{sec:conclusions-and-future-work-trabajo-futuro}.


\section{Trabajo futuro}
\label{sec:conclusions-and-future-work-trabajo-futuro}

\subsection{Sobre \CIP-133}
\label{subsec:conclusions-and-future-work-sobre-cip-133}

Una línea de trabajo futuro particularmente relevante para este proyecto es la evolución del soporte de \MSM en la capa de lenguajes y herramientas de \Cardano. En la implementación actual, el principal cuello de botella del verificador \STM no estuvo dado por la corrección criptográfica del argumento, sino por el costo de materializar \Onchain una gran \MSM sobre puntos de \BLS. Esa operación domina buena parte del costo lineal del verificador emitido en Aiken y fue una de las razones centrales por las que el flujo final debió partir la verificación en dos transacciones consecutivas.

\CIP-133 propone precisamente incorporar una \textit{builtin} nativa para \MSM sobre \BLS \cite{CIP0133, PlutusMSMBenchRepo}. Este punto merece una precisión importante: desde la perspectiva del protocolo, esa dirección ya no es meramente especulativa. El soporte para estas primitivas ya forma parte de la evolución reciente de Plutus en \Cardano, pero al momento de escribir esta tesis todavía no se encuentra expuesto de manera estable en Aiken. En particular, el trabajo para incorporar estas \textit{builtins} en el compilador y en la tipificación del lenguaje aparece reflejado en el \textit{pull request} \texttt{aiken-lang/aiken\#1349}, que al 1 de julio de 2026 seguía abierto y proponía justamente agregar soporte para las operaciones de \MSM en $\GOne$ y $\GTwo$.

Para este \zkBridge, la consecuencia potencial no sería solamente una optimización incremental. Si Aiken pasa a exponer de manera usable las primitivas asociadas a \CIP-133, una fracción importante del trabajo que hoy se expresa como cadenas largas de \texttt{scaleG1}, \texttt{addG1}, descompresiones y \textit{overhead} de la máquina CEK podría colapsar en una operación primitiva mucho más compacta. Eso volvería plausible verificar un certificado de \Mithril dentro de una sola transacción, sin necesidad de introducir el \textit{split} \texttt{phase1}/\texttt{phase2} ni de transportar un estado intermedio entre ambas fases.

La consecuencia arquitectónica sería importante. Desaparecería una parte de la complejidad hoy dedicada a los \ProofReceipt, al \textit{binding} entre fases y a la coordinación secuencial de dos transiciones para autenticar un mismo certificado. Incluso si siguieran siendo necesarios algunos \ReferenceScripts por motivos de tamaño, el diseño global del bridge podría simplificarse de manera apreciable. Por eso, una continuación natural de este trabajo sería reevaluar toda la arquitectura a la luz de \CIP-133 una vez que ese soporte quede disponible de extremo a extremo en la superficie de programación efectivamente usada por el prototipo.

\subsection{Sobre `zk-bridge-operator`}
\label{subsec:conclusions-and-future-work-sobre-zk-bridge-operator}

El comando \texttt{zk-bridge-operator tx prove <transaction-hash>} genera hoy un \textit{bundle} con dos pruebas distintas para el mismo \texttt{tx\_id} canónico de \Cardano. La primera, \texttt{snapshot\_membership}, sigue la ruta \textit{legacy} modelada en la Subsección~\ref{subsec:tx-snapshot-membership-modelo-matematico-del-transactionsnapshot-pythagoras-de-mithril} y prueba que la transacción pertenece al snapshot certificado. La segunda, \TxSetUpdateCircuit, produce la prueba de actualización del \SMT local del bridge según la relación de la Subsección~\ref{subsec:tx-set-update-la-relacion-np-de-txsetupdatecircuit}.

En el estado actual, el operador no construye aún un \textit{witness} completamente general para la segunda prueba. La función \texttt{single\_insert\_empty\_tree\_witness} arma un caso canónico de inserción sobre árbol vacío: parte de la raíz del \SMT completamente vacío, usa el camino inducido por \texttt{tx\_id} y calcula la raíz resultante al cambiar la hoja final de $0$ a $1$. Esto es suficiente para la prueba de concepto actual y no altera la semántica matemática del circuito, pero claramente estrecha el conjunto de instancias que el operador sabe producir hoy.

La consecuencia conceptual es nítida. \TxSetUpdateCircuit ya prueba la relación general \TxSetUpdateRelation: cualquier par $(r_{\mathrm{in}}, r_{\mathrm{out}})$ compatible con un mismo camino y un mismo vector de hermanos entra en su lenguaje. Lo que permanece como trabajo futuro no es generalizar el circuito, sino generalizar la capa \Offchain del operador para que construya \textit{witnesses} de actualización sobre un \SMT arbitrario mantenido a lo largo de ejecuciones sucesivas, y no sólo sobre el caso canónico de primer inserto.

Esa generalización debe incluir también la operación sostenida del servicio. El operador deberá persistir o reconstruir el estado auxiliar del \SMT, reintentar operaciones frente a cambios concurrentes del \UTxO canónico y recuperarse después de interrupciones sin producir bundles incompatibles con la raíz \Onchain vigente. Como la seguridad no exige aceptar datos por confianza en un operador particular, una evolución deseable es permitir que más de un \textit{relayer/prover} independiente construya la misma evidencia. Esto reduce la dependencia de un único proceso para la \textit{liveness}, aunque no elimina la necesidad de incentivos y mecanismos de coordinación para que alguno de ellos mantenga actualizado el protocolo.

\subsection{Sobre `preview testnet` de Cardano}
\label{subsec:conclusions-and-future-work-sobre-preview-testnet-de-cardano}

Otra continuación natural de este trabajo sería agregar soporte \Onchain real sobre la \textit{preview testnet} de \Cardano. El prototipo actual ya recorre de punta a punta el flujo transaccional del \zkBridge en un entorno local reproducible, apoyado en \Dolos, \TxThree, \Trix y \CShell. Ese recorrido, hoy orquestado por el flujo integrado del repositorio, fue suficiente para validar la lógica del protocolo, sus presupuestos de ejecución y la coherencia entre pruebas, certificados y \UTxOs. Sin embargo, todavía queda el paso de ejecutar el mismo flujo sobre una red pública de pruebas, con interacción real con la infraestructura y con las convenciones operativas de \Cardano.

En términos prácticos, esto implica integrar el bridge con una \textit{wallet} de usuario, por ejemplo \texttt{Lace}, para firmar y enviar las transacciones del flujo desde una interfaz más cercana al uso real. Ese paso permitiría validar aspectos que en el entorno local quedan parcialmente abstraídos: selección y manejo de \UTxOs desde una billetera estándar, provisión de colateral, publicación y reutilización de \ReferenceScripts en una red compartida, latencias y confirmaciones reales, y la fricción concreta de llevar el flujo desde una herramienta de desarrollo hacia un uso más cercano al de un operador o usuario final.

Desde el punto de vista de la arquitectura, esta migración no debería entenderse sólo como una etapa de despliegue, sino como una validación adicional del diseño. Si el flujo actual logra sostenerse sobre \textit{preview testnet} con firmas reales de una wallet como \texttt{Lace}, entonces la prueba de concepto dejaría de estar limitada a una demostración controlada en \textit{devnet} y pasaría a aproximarse mucho más al escenario operativo que tendría un \zkBridge usable sobre \Cardano. También permitiría estudiar con mayor seriedad la experiencia de uso, la observación de confirmaciones, la recuperación frente a transacciones rechazadas y la operación continuada del bridge, cuestiones que en esta tesis quedaron deliberadamente subordinadas a la viabilidad criptográfica y al encaje \Onchain del protocolo.

\subsection{Política de finalidad y tratamiento de \textit{rollbacks}}
\label{subsec:conclusions-and-future-work-finalidad-y-rollbacks}

El circuito de pertenencia demuestra que una \LockTx forma parte de un \TransactionSnapshot autenticado por \Mithril, pero esa propiedad no define por sí sola cuándo el bridge debe considerar suficientemente estable el estado de la cadena origen. Una versión productiva deberá incorporar una política explícita de finalidad que determine qué certificados y snapshots son admisibles para autorizar \MintTx. Esta política debe distinguir la validez criptográfica de un certificado de su vigencia operativa: un certificado correctamente firmado puede seguir siendo antiguo respecto del estado más reciente conocido por el bridge.

En particular, será necesario fijar criterios de frescura y orden. El protocolo deberá decidir qué época y qué punto certificado de la cadena resultan aceptables, cómo se relaciona el \TransactionSnapshot presentado con el \StakeDistributionUTxO vigente y qué información monotónica debe persistirse para rechazar regresiones hacia estados certificados anteriores. También deberá especificar si una transacción de origen puede reclamarse a partir de cualquier snapshot posterior que la contenga o si se exige una ventana determinada de certificación y confirmaciones.

Finalmente, la política deberá establecer cómo proceder ante un \textit{rollback} observado antes o después de un minting. Antes de aceptar \MintTx, puede imponerse una espera o una condición de estabilidad apoyada en los artefactos de \Mithril y en el punto certificado de la cadena. Después de un minting ya finalizado en la cadena destino, una reorganización de la cadena origen no puede deshacerse automáticamente mediante la lógica unidireccional actual; por lo tanto, deben definirse límites de riesgo, mecanismos de pausa o recuperación y, eventualmente, reglas de gobernanza. Formalizar y probar estas decisiones permitiría extender el modelo de seguridad de la Sección~\ref{sec:design-decisions-modelo-de-seguridad} con una propiedad temporal explícita, en lugar de tratar la finalidad como un supuesto implícito.

\subsection{Sobre \texorpdfstring{\GenesisSignature}{GenesisSignature} en la era \texorpdfstring{\Lagrange}{Lagrange}}
\label{subsec:conclusions-and-future-work-sobre-genesissignature-en-la-era-lagrange}

Un último frente de trabajo futuro queda asociado al caso particular del certificado génesis de \Mithril en la transición hacia la era \Lagrange. En el flujo desarrollado en esta tesis, el \GenesisCertificate se usa para inicializar el \StakeDistributionUTxO y se autentica directamente dentro del validador de mint mediante una \GenesisSignature verificada contra la \GenesisVerificationKey. Esa decisión fue deliberada: el certificado génesis no representa una actualización recurrente de \STM, por lo que no conviene forzarlo a pasar por el mecanismo \texttt{phase1}/\texttt{phase2} ni generar un \ProofReceipt sólo para el arranque del estado.

Sin embargo, la ruta \Lagrange introduce una dificultad criptográfica adicional. Mientras el caso más simple se reduce a una firma Ed25519 verificable con primitivas disponibles en Aiken, el formato \textit{SNARK-friendly} de \Mithril incorpora una \StandardSchnorrSignature sobre \Jubjub. Verificar esa firma directamente \Onchain no es una opción práctica en el corto plazo: \Jubjub no es una curva nativa de \Cardano y una implementación pura en Aiken de aritmética de curva, multiplicación escalar y recomputación del desafío consumiría demasiadas unidades de ejecución. Por lo tanto, una extensión natural del sistema es mover esa verificación a un circuito auxiliar y dejar que el validador de \StakeDistribution verifique sólo una prueba sucinta.

Concretamente, ese trabajo consiste en definir un circuito \Groth sobre \BLS que capture la relación de verificación Schnorr usada por el génesis de \Lagrange: reconstruir el punto $R' = rG + c\SchnorrVK$, recomputar el desafío con la variante de \PoseidonHash/\MidnightPoseidonHash usada por \Mithril sobre el \textit{domain separation tag}, la clave, el punto reconstruido y el mensaje base, y exigir que ese desafío coincida con el de la firma. Además, el circuito debe fijar las condiciones de pertenencia al subgrupo correcto y conectar sus entradas públicas con la codificación real del certificado, no sólo con valores algebraicos ya parseados.

Desde el punto de vista del bridge, la integración esperada mantendría la misma filosofía del flujo final: \StakeDistributionGenesisTxTemplate seguiría siendo la transacción de bootstrap y no reaparecería un dominio \texttt{stake\_distribution\_genesis} dentro de las dos fases de verificación \STM. La diferencia sería que el redeemer del caso génesis podría transportar, además de la firma Ed25519, una prueba \Groth de la firma Schnorr sobre \Jubjub. El validador combinaría entonces dos chequeos: la autenticación nativa con \GenesisVerificationKey y la verificación del argumento sucinto que ata la componente \Lagrange del certificado al mismo mensaje génesis.

La transición requiere además una política de versionado para el material de verificación. La versión del certificado, la relación aritmetizada, la clave de verificación y el script Aiken que la consume deben tratarse como una unidad de despliegue: cambiar solamente uno de esos elementos puede volver incompatibles pruebas y estados previamente publicados. Una implementación productiva deberá identificar explícitamente cada versión, autenticar las nuevas claves de verificación y definir cómo migra el \StakeDistributionUTxO vigente sin abrir una segunda raíz de confianza ni dejar receipts de una versión antigua utilizables bajo la nueva.

Aunque una versión posterior del prototipo ya avanzó en esta dirección, este trabajo queda correctamente clasificado como fuera del alcance original de la tesis.
